\originalchapter{Smith 标准形与模分解}
\originalsection{关系矩阵的化简}
给关系矩阵换行对应改变自由模 $R^s$ 的基, 换列对应改变关系的生成列, 两者都不改变商模的同构类型. 允许的初等变换包括交换行列、某一行列乘单位、给一行列加另一行列的倍数, 它们均可逆.
\begin{theorem}[Smith 标准形]\label{thm:smith}
PID $R$ 上的矩阵 $A$ 可经可逆行列变换化为对角矩阵, 非零对角元 $d_1,\ldots,d_r$ 满足 $d_1\mid d_2\mid\cdots\mid d_r$.
\end{theorem}
\begin{proof}
零矩阵已经是所需对角形. 以下非零情形先说明两元素的消元, 其中 $(a,b)\ne(0,0)$. 对 $a,b$ 取 gcd $d=ua+vb$, 写 $a=da'$, $b=db'$, 则 $ua'+vb'=1$. 矩阵
\[
\begin{pmatrix}u&v\\-b'&a'\end{pmatrix}
\]
行列式为1, 将列向量 $(a,b)^T$ 变成 $(d,0)^T$. 同理可作列变换.

把一个非零元素移到左上角. 若首列某元素不能被左上角元素整除, 用上述 gcd 变换把左上角改成严格较大的主理想生成元; 首行也同样处理. 这种严格增大不可能无限继续, 因 PID 的理想升链停止. 因而终可使左上角整除首行首列所有项, 再用减去倍数清零, 得 $\operatorname{diag}(d,B)$.

若 $B$ 中仍有项 $b$ 不被 $d$ 整除, 将该项所在行加到首行, 首行在该列出现 $b$, 再用 gcd 列消元, 左上角主理想又严格增大. 升链停止保证有限次后 $d$ 整除 $B$ 中每项. 对 $B$ 归纳. 在其内部的行列变换仍使每项被 $d$ 整除, 故其第一个非零对角元也被 $d$ 整除. 归纳给出整除链.
\end{proof}
\originalsection{有限生成模结构定理}
\begin{theorem}\label{thm:module}
PID 上任何有限生成模同构于
\[
R^t\oplus R/(d_1)\oplus\cdots\oplus R/(d_r),\qquad d_1\mid\cdots\mid d_r,
\]
其中 $d_i$ 为非零非单位. 自由秩 $t$ 和各 $d_i$ 的相伴类唯一. 等价地, 挠部分可按素元幂分解为初等因子.
\end{theorem}
\begin{proof}
由上一章, $M$ 有有限关系矩阵. Smith 变换给出直和商: 非零对角元 $d$ 给 $R/(d)$, 零行给自由因子, 单位对角元给零模, 可删去.

对唯一性, 先把模延标量到分式域, 每个非零 $d_i$ 在其中可逆, 所以维数为 $t$, 自由秩唯一. 对素元 $p$, 将挠模中的 $p$ 主部分记为 $T_p=\{m:p^k m=0\text{ 对某 }k\}$. 对互素因子的中国剩余定理说明挠模是这些主部分的直和. 在 $T_p$ 上, $p^{k-1}T_p/p^kT_p$ 是域 $R/(p)$ 上的向量空间. 一个因子 $R/(p^a)$ 在该维数中贡献1当且仅当 $a\ge k$, 记该维数为 $b_k$, 则次数恰为 $k$ 的初等因子有 $b_k-b_{k+1}$ 个, 所以这些维数恢复每个素元幂指数的多重集. 最后将每种素元的指数按非降序排列令 $r$ 为各素元初等因子数的最大值 (挠模为0时取 $r=0$), 并在左侧补零至长度 $r$, 原整除链中首个非单位因子必含某素元, 该素元因整除链而出现在全部因子中, 因而原非单位因子数就是此最大值 $r$. 逐位置相乘, 恢复整除链的 $d_i$; 这也证明初等因子与不变因子之间转换唯一.
\end{proof}
这里延标量可具体理解为形式分数 $m/s$, $s\ne0$, 以 $m/s=m'/s'$ 当某非零 $u$ 杀掉 $s'm-sm'$ 来识别; 具体为 $m/s+m'/s'=(s'm+sm')/(ss')$, $(a/b)(m/s)=am/(bs)$. 等价式传递性通过相乘非零分母及两个原见证因子得到新见证; 改变代表元后, 加法或标量乘法两边的交叉乘积之差是原零差的标量线性组合, 再乘原见证因子的乘积即可杀掉, 故运算良定义. 向量空间各公理由分配律和分式运算直接继承. 模同构 $f$ 延拓为 $m/s\mapsto f(m)/s$, 其逆由 $f^{-1}$ 延拓, 因而保留延标量维数. 对直和分量逐坐标取分数即可得到相应直和. 其核恰为挠元, 对自由因子得到 $K^t$, 因而上述自由秩论证不需要预设张量积理论.
\begin{example}分解由矩阵 $\begin{pmatrix}4&6\\2&8\end{pmatrix}$ 给出的 $\Z$ 模.\end{example}
\begin{solution}
所有元素的 gcd 是2, 行列式为20. Smith 形第一个因子为2, 第二个为10: 可先交换行, 将第二行减第一行2倍得 $(0,-10)$, 首行是 $(2,8)$, 再将第二列减第一列4倍清掉8, 最后改变符号. 商为 $\Z/2\Z\oplus\Z/10\Z$, 阶20.
\end{solution}
\originalsection{线性算子的代数解释}
给域 $F$ 上有限维空间 $V$ 和线性算子 $T$, 定义 $f(x)v=f(T)v$, 使 $V$ 成 $F[x]$ 模. 它全为挠模: 任意 $v$, 向量 $v,Tv,\ldots,T^nv$ 线性相关, 所以某非零多项式杀掉 $v$. 结构定理给出循环因子 $F[x]/(f_i)$, $f_i\mid f_{i+1}$. 在基 $1,x,\ldots,x^{d-1}$ 中, 乘 $x$ 的矩阵是 $f_i$ 的伴随矩阵, 得到有理标准形. 若多项式分裂, 初等因子 $F[x]/((x-\lambda)^a)$ 在基 $1,(x-\lambda),\ldots$ 中给出一个 Jordan 块. 此处是模分解的应用, 不展开数值计算或稳定性问题.
\originalsection{带解答练习}
\begin{exercise}列出阶72的所有交换群同构类型.\end{exercise}
\begin{solution}
$72=2^3 3^2$. 2主部分为 $C_8,C_4\times C_2,C_2^3$; 3主部分为 $C_9,C_3^2$. 任取两类直积, 共6种. 唯一性由初等因子保证, 所以不会重复.
\end{solution}
\begin{exercise}若 $T^2=0$ 且 $\dim V=5$, $\operatorname{rank}T=2$, 求 Jordan 块.\end{exercise}
\begin{solution}
仅能有大小1或2的零特征值块, 每个大小2块秩为1. 因而有两个2块, 剩余一个1块. 无须假设底域代数闭, 因最小多项式已经分裂为 $x^2$ 的因子.
\end{solution}
